% =============================================================
%  Section II — System Model and Problem Formulation
%  File: sections/04_system_model.tex
%  Note: \input'd from main.tex — do NOT include preamble here
% =============================================================

\section{System Model and Problem Formulation}
\label{sec:system}

We consider the same monostatic MIMO-OFDM ISAC framework as in~\cite{ldsa}, but impose an additional \emph{RBG-granularity constraint} that reflects the actual scheduling structure mandated by the 5G~NR physical layer~\cite{3gpp.38.214}.
Throughout this section, we first recap the notation and key results from~\cite{ldsa} (Sections~\ref{subsec:signal}--\ref{subsec:crb_recap}), then introduce the RBG structure (Section~\ref{subsec:rbg_structure}) and formulate the new optimisation problem (Section~\ref{subsec:problem}).

% -------------------------------------------------------------
\subsection{Signal Model}
\label{subsec:signal}
% -------------------------------------------------------------

A monostatic ISAC base station (BS) is equipped with a uniform linear array (ULA) of $N_T$ transmit and $N_R$ receive antennas with half-wavelength spacing $d = \lambda/2$.
The system employs an OFDM waveform with $N$ subcarriers at subcarrier spacing $\Delta f$ and carrier frequency $f_c$.
A coherent processing interval (CPI) comprises $L$ OFDM symbols with symbol duration~$T_s$.
The BS simultaneously performs multi-target radar sensing and serves $U$ downlink OFDMA users.

\subsubsection{Communication}
On each communication subcarrier~$n$ (i.e., $w_n = 0$), the BS transmits to the designated user $u_n \in \{1,\ldots,U\}$ with water-filling power $p_n$.
The per-subcarrier achievable rate is
\begin{equation}
  r_n = \log_2\!\left(1 + \frac{p_n |h_{n,u_n}|^2}{\sigma_c^2}\right),
  \label{eq:rate_sub}
\end{equation}
where $h_{n,u_n}$ is the frequency-domain channel coefficient from the BS to user~$u_n$ on subcarrier~$n$, and $\sigma_c^2$ is the noise power at the communication receiver.
The aggregate rate delivered to user~$u$ under a binary allocation vector $\mathbf{w} = [w_0, \ldots, w_{N-1}]^\top$ is
\begin{equation}
  R_u(\mathbf{w}) = \sum_{\substack{n:\, u_n = u}} (1 - w_n)\, r_n.
  \label{eq:rate_user}
\end{equation}

\subsubsection{Sensing}
After matched filtering, the baseband echo from a point target at range $R$, velocity $v$, and angle-of-arrival (AoA) $\theta$ is observed at receive antenna $r$, subcarrier $n$, and OFDM symbol $l$ as~\cite{sturm2011}
\begin{equation}
  y[r,n,l] = \beta \, e^{j2\pi \frac{d}{\lambda}\sin\theta \cdot r}
              \, e^{-j2\pi n \Delta f \tau}
              \, e^{j2\pi f_D l T_s}
              + z[r,n,l],
  \label{eq:echo}
\end{equation}
where $\tau = 2R/c$ is the round-trip delay, $f_D = 2v f_c / c$ is the Doppler frequency, $\beta \in \mathbb{C}$ is the complex target amplitude, and $z[r,n,l] \sim \mathcal{CN}(0, \sigma_z^2)$ is additive white Gaussian noise.
Only sensing subcarriers ($w_n = 1$) contribute observations.

% -------------------------------------------------------------
\subsection{CRB Analysis (Recap from~\cite{ldsa})}
\label{subsec:crb_recap}
% -------------------------------------------------------------

We briefly recall the key result of~\cite{ldsa} that motivates the objective function.
The structural parameter vector is $\boldsymbol{\eta} = [\tau, f_D, \theta]^\top$ (after concentrating out $\beta$).
Define the moment sums of the selected subcarrier indices:
\begin{equation}
  S_k(\mathbf{w}) = \sum_{n=0}^{N-1} w_n \, n^k, \quad k = 0, 1, 2,
  \label{eq:moments}
\end{equation}
so that $S_0 = K$ (the sensing budget), $\bar{n} = S_1 / K$ is the index mean, and the index variance is
\begin{equation}
  V(\mathbf{w}) = \frac{1}{K} \sum_{n=0}^{N-1} w_n (n - \bar{n})^2
               = \frac{S_2}{K} - \left(\frac{S_1}{K}\right)^{\!2}.
  \label{eq:variance}
\end{equation}

As shown in~\cite{ldsa}, by adopting centred subcarrier coordinates $\tilde{n} = n - \bar{n}$, the Fisher Information Matrix (FIM) block-diagonalises into a delay block and a Doppler--angle block:
\begin{equation}
  \FIM^{(c)}(\mathbf{w}) = \gamma \, \blkdiag\!\Big(
    N_R L \, \alpha_f^2 K \, V(\mathbf{w}),\;
    \mathbf{J}_{23}
  \Big),
  \label{eq:fim_block}
\end{equation}
where $\gamma = 2|\beta|^2 / \sigma_z^2$, $\alpha_f = 2\pi \Delta f$, and $\mathbf{J}_{23}$ is the $2\times 2$ Doppler--angle block that depends on $K$ but \emph{not} on the selection pattern.
The range Cram\'{e}r--Rao bound (CRB) follows as
\begin{equation}
  \CRB(R;\, \mathbf{w}) = \frac{c^2}{16\pi^2 \Delta f^2 \cdot \gamma \cdot N_R \cdot L \cdot K \cdot V(\mathbf{w})}.
  \label{eq:crb_range}
\end{equation}
Since all parameters except $V(\mathbf{w})$ are fixed by system design, \emph{minimising} $\CRB(R)$ is equivalent to \emph{maximising} the index variance $V(\mathbf{w})$.

% -------------------------------------------------------------
\subsection{RBG Structure in 5G NR}
\label{subsec:rbg_structure}
% -------------------------------------------------------------

In the existing per-subcarrier formulation~\cite{ldsa}, $N$ independent binary variables $\{w_n\}_{n=0}^{N-1}$ are optimised.
However, the 5G NR standard specifies that downlink frequency-domain resource allocation is performed at the granularity of \emph{resource block groups} (RBGs)~\cite{3gpp.38.214}.

\begin{definition}[RBG Structure]
\label{def:rbg}
Each physical resource block (PRB) comprises 12 consecutive subcarriers.
An RBG consists of $P$ consecutive PRBs, where $P \in \{2, 4, 8, 16\}$ is determined by the bandwidth part (BWP) size according to TS~38.214, Table~5.1.2.2.1-1~\cite{3gpp.38.214}.
The total number of RBGs is
\begin{equation}
  G = \left\lfloor \frac{N}{12P} \right\rfloor.
  \label{eq:num_rbg}
\end{equation}
\end{definition}

Let $M = 12P$ denote the number of subcarriers per RBG.
The $g$-th RBG ($g = 1, \ldots, G$) occupies the contiguous subcarrier index set
\begin{equation}
  \Gcal_g = \big\{(g-1)M,\; (g-1)M + 1,\; \ldots,\; gM - 1\big\}.
  \label{eq:rbg_set}
\end{equation}
Since all $M$ subcarriers within an RBG must be jointly allocated to either sensing or communication, we introduce the \emph{RBG-level binary decision variable}:
\begin{equation}
  \tilde{w}_g \in \{0, 1\}, \quad g = 1, \ldots, G,
  \label{eq:rbg_var}
\end{equation}
where $\tilde{w}_g = 1$ assigns all subcarriers in $\Gcal_g$ to sensing, and $\tilde{w}_g = 0$ assigns them to communication.
The per-subcarrier allocation is then determined as $w_n = \tilde{w}_g$ for all $n \in \Gcal_g$.

\begin{remark}[Reduction to per-subcarrier case]
\label{rmk:reduction}
Setting $P = 1/12$ (i.e., $M = 1$) recovers the per-subcarrier formulation in~\cite{ldsa}.
The RBG constraint is therefore a strict generalisation that subsumes the idealised model as a limiting case.
\end{remark}

% -------------------------------------------------------------
\subsection{RBG-Level Aggregates}
\label{subsec:rbg_aggregates}
% -------------------------------------------------------------

Under the RBG constraint, the moment sums~\eqref{eq:moments} and the variance~\eqref{eq:variance} can be expressed entirely in terms of the $G$-dimensional vector $\tilde{\mathbf{w}} = [\tilde{w}_1, \ldots, \tilde{w}_G]^\top$ through pre-computable per-RBG aggregates.

\begin{definition}[Per-RBG aggregates]
\label{def:aggregates}
For each RBG $g = 1, \ldots, G$, define:
\begin{subequations}\label{eq:rbg_agg}
\begin{align}
  \mu_g &= \sum_{n \in \Gcal_g} n
         = (g-1)M^2 + \tfrac{M(M-1)}{2},
  \label{eq:mu_g} \\
  \nu_g &= \sum_{n \in \Gcal_g} n^2
         = \sum_{k=0}^{M-1} \big((g-1)M + k\big)^2,
  \label{eq:nu_g} \\
  \bar{r}_{g,u} &= \sum_{\substack{n \in \Gcal_g \\ u_n = u}} r_n, \quad \forall\, u \in \Ucal,
  \label{eq:rate_rbg}
\end{align}
\end{subequations}
where $\mu_g$ and $\nu_g$ are the first and second moment sums of subcarrier indices within RBG~$g$, and $\bar{r}_{g,u}$ is the aggregate achievable rate on RBG~$g$ for user~$u$.
Both $\mu_g$ and $\nu_g$ are computable in $O(1)$ per RBG via standard summation formulae.
\end{definition}

In terms of these aggregates, the moment sums become:
\begin{equation}
  S_k(\tilde{\mathbf{w}}) = \sum_{g=1}^{G} \tilde{w}_g \sum_{n \in \Gcal_g} n^k, \quad k = 0, 1, 2.
  \label{eq:moments_rbg}
\end{equation}
Specifically, $S_0 = K_G M$ where $K_G = \sum_{g=1}^{G} \tilde{w}_g$ is the number of sensing RBGs (total sensing subcarrier count $K = K_G M$), $S_1 = \sum_g \tilde{w}_g \mu_g$, and $S_2 = \sum_g \tilde{w}_g \nu_g$.

The index variance is then
\begin{equation}
  V(\tilde{\mathbf{w}}) = \frac{\sum_{g=1}^{G} \tilde{w}_g \, \nu_g}{K_G M}
                         - \left(\frac{\sum_{g=1}^{G} \tilde{w}_g \, \mu_g}{K_G M}\right)^{\!2},
  \label{eq:var_rbg}
\end{equation}
and the per-user aggregate rate becomes
\begin{equation}
  R_u(\tilde{\mathbf{w}}) = \sum_{g=1}^{G} (1 - \tilde{w}_g)\, \bar{r}_{g,u}, \quad \forall\, u \in \Ucal.
  \label{eq:rate_rbg_total}
\end{equation}

% -------------------------------------------------------------
\subsection{Intra-/Inter-RBG Variance Decomposition}
\label{subsec:variance_decomp}
% -------------------------------------------------------------

A key structural property of the RBG formulation is that the total index variance admits a clean decomposition into \emph{inter-RBG} and \emph{intra-RBG} components, analogous to the analysis of variance (ANOVA).

\begin{proposition}[Variance decomposition]
\label{prop:var_decomp}
Let $c_g = (g - 1)M + (M-1)/2$ denote the centroid (mean index) of RBG~$g$, and let
\begin{equation}
  V_{\mathrm{intra}} = \frac{M^2 - 1}{12}
  \label{eq:v_intra}
\end{equation}
be the common intra-RBG variance (identical for all RBGs since each contains $M$ consecutive integers).
Then
\begin{equation}
  V(\tilde{\mathbf{w}}) = \underbrace{\frac{1}{K_G} \sum_{g=1}^{G} \tilde{w}_g \big(c_g - \bar{c}\big)^2}_{V_{\mathrm{inter}}(\tilde{\mathbf{w}})}
                        + \; V_{\mathrm{intra}},
  \label{eq:anova}
\end{equation}
where $\bar{c} = \frac{1}{K_G}\sum_{g=1}^{G} \tilde{w}_g \, c_g$ is the weighted centroid mean.
\end{proposition}

\begin{proof}
For any selected subcarrier $n \in \Gcal_g$ with $\tilde{w}_g = 1$, write $n = c_g + \delta_n$ where $\delta_n = n - c_g$ is the offset from the RBG centroid.
Since each RBG contains $M$ consecutive integers, $\sum_{n \in \Gcal_g} \delta_n = 0$ for every $g$, and the overall subcarrier mean satisfies $\bar{n} = \bar{c}$.
Expanding $(n - \bar{n})^2 = (c_g - \bar{c})^2 + 2(c_g - \bar{c})\delta_n + \delta_n^2$ and summing over all selected subcarriers:
\begin{equation}
  \sum_{g:\,\tilde{w}_g=1} \sum_{n \in \Gcal_g} (n - \bar{n})^2
  = M \sum_{g:\,\tilde{w}_g=1} (c_g - \bar{c})^2
    + K_G \cdot \frac{M(M^2 - 1)}{12},
  \label{eq:decomp_expand}
\end{equation}
where the cross-term vanishes due to $\sum_{n \in \Gcal_g} \delta_n = 0$.
Dividing both sides by $K = K_G M$ yields~\eqref{eq:anova}.
\end{proof}

\begin{remark}[Implication for optimisation]
\label{rmk:opt_implication}
Since $V_{\mathrm{intra}}$ is a constant determined solely by $M = 12P$, maximising $V(\tilde{\mathbf{w}})$ reduces to maximising $V_{\mathrm{inter}}(\tilde{\mathbf{w}})$, the \emph{variance of the RBG centroids}.
This is the RBG-level analogue of the per-subcarrier result in~\cite{ldsa}:
instead of dispersing individual subcarriers, the optimiser now disperses RBG \emph{blocks}.
The irreducible sensing loss from the grouping constraint is captured entirely by $V_{\mathrm{intra}}$, which grows quadratically with the RBG size~$M$.
\end{remark}

% -------------------------------------------------------------
\subsection{Problem Formulation}
\label{subsec:problem}
% -------------------------------------------------------------

Using the variance decomposition of Proposition~\ref{prop:var_decomp}, we formulate the RBG-constrained ISAC subcarrier allocation as:

\begin{subequations}\label{eq:P0_RBG}
\begin{align}
  (\text{P0-RBG}): \quad \max_{\tilde{\mathbf{w}}} \quad
    & V_{\mathrm{inter}}(\tilde{\mathbf{w}})
    = \frac{1}{K_G}\sum_{g=1}^{G} \tilde{w}_g\, c_g^2
      - \left(\frac{1}{K_G}\sum_{g=1}^{G} \tilde{w}_g\, c_g\right)^{\!2}
    \label{eq:P0_obj} \\
  \text{s.t.} \quad
    & \text{C1:} \quad \sum_{g=1}^{G} \tilde{w}_g = K_G,
    \label{eq:P0_C1} \\
    & \text{C2:} \quad R_u(\tilde{\mathbf{w}}) \geq \Rmin_u, \quad \forall\, u \in \Ucal,
    \label{eq:P0_C2} \\
    & \text{C3:} \quad \tilde{w}_g \in \{0, 1\}, \quad \forall\, g = 1, \ldots, G.
    \label{eq:P0_C3}
\end{align}
\end{subequations}

Here, $K_G$ is the prescribed number of sensing RBGs (equivalently, $K = K_G M$ sensing subcarriers).
Constraint C1 fixes the sensing budget at the RBG level.
Constraint C2 enforces per-user quality-of-service (QoS), where the per-user rate is given by~\eqref{eq:rate_rbg_total}.
Constraint C3 is the integrality requirement imposed by the NR scheduling granularity.

\begin{remark}[Relationship to (P0) of~\cite{ldsa}]
\label{rmk:relationship}
Problem~(P0-RBG) has $G$ binary variables versus the $N$ variables in~\cite{ldsa}, with $G = N/(12P) \ll N$.
The objective~\eqref{eq:P0_obj} has an identical mathematical structure---a ratio of linear and quadratic terms in~$\tilde{\mathbf{w}}$---but operates on the RBG centroid indices $\{c_g\}$ rather than individual subcarrier indices.
\end{remark}

\subsubsection{Convex relaxation}
We relax C3 to $\tilde{w}_g \in [0, 1]$:
\begin{subequations}\label{eq:P1_RBG}
\begin{align}
  (\text{P1-RBG}): \quad \max_{\tilde{\mathbf{w}} \in [0,1]^G} \quad
    & f(\tilde{\mathbf{w}}) = \frac{1}{K_G}\sum_{g} \tilde{w}_g\, c_g^2
      - \frac{1}{K_G^2}\!\left(\sum_{g} \tilde{w}_g\, c_g\right)^{\!2}
    \label{eq:P1_obj} \\
  \text{s.t.} \quad
    & \text{C1, C2 in~\eqref{eq:P0_RBG}}.
    \label{eq:P1_constraints}
\end{align}
\end{subequations}

\begin{proposition}[Concavity of (P1-RBG)]
\label{prop:concavity}
The objective $f(\tilde{\mathbf{w}})$ in~\eqref{eq:P1_obj} is concave on the feasible set $\{\tilde{\mathbf{w}} : \sum_g \tilde{w}_g = K_G,\; \tilde{\mathbf{w}} \in [0,1]^G\}$.
Consequently, (P1-RBG) is a convex programme.
\end{proposition}

\begin{proof}
On the affine subspace defined by C1, $K_G = \sum_g \tilde{w}_g$ is constant.
The first term $\frac{1}{K_G}\sum_g \tilde{w}_g c_g^2$ is linear in $\tilde{\mathbf{w}}$.
The second term can be written as $-\frac{1}{K_G^2}\, \tilde{\mathbf{w}}^\top \mathbf{c}\mathbf{c}^\top \tilde{\mathbf{w}}$ where $\mathbf{c} = [c_1, \ldots, c_G]^\top$.
Since $\mathbf{c}\mathbf{c}^\top$ is positive semidefinite (rank one), this is a concave quadratic.
The sum of a linear and a concave function is concave.
All constraints are affine or box constraints, hence the feasible set is convex.
\end{proof}

\begin{remark}[Strong duality]
\label{rmk:slater}
Strong duality holds for (P1-RBG) provided Slater's condition is satisfied, i.e., $K_G < G$ and $\sum_g \bar{r}_{g,u} > \Rmin_u$ for all $u \in \Ucal$.
This enables the Lagrangian dual approach developed in Section~\ref{sec:method}.
\end{remark}

\subsubsection{Structural comparison}
Table~\ref{tab:comparison} summarises the key differences between the per-subcarrier and RBG-level formulations.

% tables/table_comparison.tex
\begin{table}[t]
  \centering
  \caption{Per-Subcarrier vs.\ RBG-Level Formulation Comparison}
  \label{tab:comparison}
  \begin{tabular}{@{}lcc@{}}
    \toprule
    \textbf{Property} & \textbf{Per-SC~\cite{ldsa}} & \textbf{RBG (this work)} \\
    \midrule
    Decision variables       & $N$ ($w_n$)       & $G = N/(12P)$ ($\tilde{w}_g$) \\
    Granularity              & 1 subcarrier      & $M = 12P$ subcarriers \\
    Index used in objective  & $n$               & $c_g$ (RBG centroid) \\
    Intra-block variance     & 0                 & $(M^2 - 1)/12$ \\
    NR DCI-compatible        & No                & Yes \\
    Typical $G/N$ ratio      & 1                 & $1/24$ to $1/192$ \\
    \bottomrule
  \end{tabular}
\end{table}

